\section{Discussion et conclusion}

\paragraph{L'approche la plus appropriée.} Prédire séparément résistance, ductilité et
ténacité, puis les comparer aux seuils d'une norme, avec une forêt aléatoire réglée par
validation croisée imbriquée groupée par article. Elle surpasse les modèles linéaires et
la réduction de dimension ; sa variante semi-supervisée ne fait pas nettement mieux, et
nous gardons la forêt supervisée, plus simple. Mais ses performances restent modestes
sur une étude inconnue ($R^2$ de 0,20 à 0,46), et le verdict tiré de ses prédictions
laisse passer la plupart des soudures mauvaises : en l'état, le modèle peut trier des
recettes candidates, pas remplacer les essais.

\paragraph{Recommandations, avec prudence.} Les tendances suivantes ressortent à la fois
des corrélations, des modèles et de la métallurgie :
\begin{itemize}
  \item les éléments d'alliage (Mo, Mn, V, Cr, Nb) augmentent la résistance mais
    réduisent la ductilité : une bonne soudure est un compromis, à régler selon l'usage ;
  \item un traitement thermique plus chaud ou plus long abaisse la résistance ;
  \item un fort apport de chaleur au soudage et les impuretés (S, P) vont avec une
    striction plus faible ; l'énergie de soudage, qui intègre courant, tension et
    vitesse, est le réglage à suivre ;
  \item la ténacité se juge à la température d'usage : une soudure satisfaisante à
    \SI{0}{\celsius} peut être fragile à \SI{-40}{\celsius}.
\end{itemize}
Ce sont des corrélations issues de données publiées, pas des causes : les aciers alliés
étant presque toujours traités, l'effet des alliages et celui du traitement ne se
séparent pas. Des essais contrôlés, menés avec des experts matériaux, seraient
nécessaires pour les confirmer.

\paragraph{Pistes.} Des données plus variées (plus d'articles et d'aciers), des modèles
par famille d'acier, un verdict prudent (marge sur le seuil, régression par quantile ou
classification apprise sur la zone du seuil), et un traitement à part des énergies
Charpy fixées à \SI{100}{\joule} et \SI{28}{\joule}, qui semblent résumer une courbe
plutôt que mesurer un essai.

\paragraph{Organisation de l'équipe.} La figure~\ref{fig:gantt} résume le partage des
tâches.

% Diagramme construit à partir de l'historique du dépôt git : dates et répartition
% à vérifier par l'équipe avant le rendu.
\begin{figure}[!htbp]
  \centering\small
  \begin{ganttchart}[
      time slot format=isodate,
      x unit=0.62cm, y unit chart=0.42cm, y unit title=0.5cm,
      title height=1, bar height=0.6, bar top shift=0.2,
      title label font=\scriptsize, bar label font=\scriptsize,
      bar/.append style={fill=NavyBlue!55, draw=NavyBlue},
      vgrid={*1{black!10}}, hgrid
    ]{2026-09-30}{2026-10-14}
    \gantttitlecalendar{day} \\
    \ganttbar{Récupération et description des données (Margaux)}{2026-09-30}{2026-10-02} \\
    \ganttbar{Exploration des données (Myriam, Salomé)}{2026-10-02}{2026-10-03} \\
    \ganttbar{Nettoyage et pré-traitement (Margaux, Salomé, Pablo)}{2026-10-02}{2026-10-08} \\
    \ganttbar{ACP (Myriam, Salomé)}{2026-10-03}{2026-10-11} \\
    \ganttbar{Protocole commun, arbres et forêts (Salomé)}{2026-10-09}{2026-10-10} \\
    \ganttbar{Régression linéaire (Pablo)}{2026-10-08}{2026-10-13} \\
    \ganttbar{Réduction de dimension (Myriam)}{2026-10-11}{2026-10-13} \\
    \ganttbar{Semi-supervisé (Margaux)}{2026-10-11}{2026-10-12} \\
    \ganttbar{Verdict ISO 2560 (Pablo)}{2026-10-13}{2026-10-13} \\
    \ganttbar{Rédaction des rapports (Salomé, Pablo)}{2026-10-03}{2026-10-14}
  \end{ganttchart}
  \caption{Diagramme de Gantt : tâches, durée et membres de l'équipe (jours de
    septembre et octobre 2026).}
  \label{fig:gantt}
\end{figure}
